\documentclass[11pt]{article}
\usepackage{amsmath,amssymb,hyperref}
\title{An egress hold can consume the data mule's last upload slot}
\author{}\date{}
\begin{document}\maketitle

\paragraph{Scope.} This is held-in I2 calibration, not a new independent
layout or a frozen stack. The topology is the already inspected
\href{https://github.com/LTU-RAI/SubTGraph/tree/21a1bcabf37c8f2f6f236ba40b916094d9903b5b/benchmark}{SubTGraph natural/01 grid},
as normalized by the parent study. Verify the actual adjacency, but do not
infer travel times, equipment, permitted edges, optical visibility, link
budgets or correlated failures from occupied cells. These are all explicit
fixture inputs. No collision or packet-delivery measurements are claimed.

\paragraph{Mission and local interface.}
An inspection in cell \(X=(11,28)\) is underway before the fault
and finishes at tick 1. The inspector normally leaves by \(X\to D=(11,27)\).
A blockage removes this edge and changes the executable continuation to
\(X\to V=(12,28)\to E=(13,28)\). The last leg uses a powered fixture
whose operator must stay at \(R=(12,29)\) until the inspector is clear at
\(E\). This dependency is a stated local hardware interlock, not a
reservation or a scheduler acknowledgment. The egress authorization
allows only these two inspector routes; the rest of the source grid is
not an alternate route search domain for this result. The inspector may
wait at \(V\) under the persistent hold while its last edge is unavailable.
In all admitted blocked cases, the hold must arm at tick 2 and physical
clearance would occur at an admitted \(c\ge4\) if it is kept. Missing the
arm cut gives no egress by the declared deadline; the inspector does not
enter under an unaccepted command. Other fallback survival is unassessed.

The mule starts at \(U=(11,29)\). It gets the local hold request and, if
the primary route to the base is unavailable, a report copy from the
adjacent inspector at the tick-1 handoff. The fixture makes the needed
status and completed exchange locally available at that cut. No shared
hidden map or post-partition live planner state is supplied. Tick labels
are elapsed from the local handoff and transport is synchronous with
locally readable slot boundaries; clock drift and establishment of those
boundaries are outside the kernel. The mule can go directly \(U\to T=(10,29)\)
to upload or instead arm \(U\to R\) at 2 and stay through clearance. From
\(R\), reaching \(T\) via \(U\) costs two ticks. Earliest possible departure
uses the local positive crossing indication, or the declared latest-clear
bound if that indication is unavailable. The timer is justified only by
the admitted persistent hold and crossing law, not silence. The report
and inspector are separate: the inspector keeps its copy, and publication
is settled only by a final acknowledgment for that report. Upload attempts
have one tick of duration, observe success only after an attempted send,
and consume one unit of the mule's attempt budget even on a failed slot.
No delayed acknowledgments, concurrent outstanding attempts, other
uploader observations, multi-packet effects or alternate upload contacts
in this fixture. A separate worker completes its local task at tick 2
even in blackout; its routes and resources are disjoint.

Nominally both robots reach authorized egress points at 2 and the base
acknowledges at 2. With only primary contact loss the mule uploads at
tick 2 and is acknowledged at 3. With only the blockage the inspector
reaches \(E\) at 4, the mule reaches \(T\) at 6, and the independent
primary upload was acknowledged at 2. With both faults, the duty to
deliver the report now depends on when the mule can leave its hold and
which upload slots remain. An independent report-first plan can deliver
at 3 while missing inspector egress; even returning immediately from
\(T\) cannot reach \(R\) by 2. This negative control does not rule out
all independent or decentralized policy families. The joint envelope
has a second, pre-settlement fault in the form of loss of a candidate
backup upload slot. Waiting for global return of the primary link is
not used.

\paragraph{Exact attempt budget for this route and observation class.}
Let each admitted history \(w\) give a clearance \(c_w\) and set \(G_w\)
of successful send slots at \(T\). Given the admitted sensor and clearance
bound, let \(a_w\) be the earliest safe arrival under the above interlock.
It is \(c_w+2\) with exact live clearance, and \(C+2\) if only the common
guaranteed latest clearance \(C\) is available. In the latter comparison
there is no other information between arm and departure. In the former,
the observed arrival derived from the hold is available to the uploader.
Under these hypotheses, earliest safe departure weakly dominates delaying:
leave as soon as permitted and wait at \(T\) to emulate any later feasible
arrival. Prior to the first upload attempt, \(a_w\) is all the source
information the policy consumes. If there is a later unsolicited status
or a way to read channel health at \(T\), this reduction need not hold.

For an arrival observation \(a\), the safe finite slot interval is
\(B(a)=\{a,\ldots,H-1\}\), where \(H\) is the acknowledgment deadline.
The least number of attempts sufficient in the worst case is
\[
 \max_{a}\ \min\{|S|:S\subseteq B(a),\
       \forall w\ (a_w=a\ \Longrightarrow\ S\cap G_w\ne\varnothing)\}.
\]
An empty feasible family for the inner minimization means infinity.
Sufficiency: issue the selected slots in order until acknowledgment.
Necessity: for each arrival observation, follow the policy on its
consecutive no-acknowledgment branch, extending arbitrarily after that
branch becomes infeasible. Truncate the branch at the budget or deadline.
All histories in the class follow it until their first successful slot,
so its at-most-budget slots must meet every \(G_w\). A randomized rule
requiring a guarantee for every seed offers no improvement. No claim
here concerns positive expected service with smaller budgets.

For \(H=8\), take \((c,G)=(4,\{6\})\) or \((5,\{7\})\). A local clearance
observation gives a different physical arrival in each case; one upload
attempt suffices. Waiting to tick 5 for guaranteed release misses the
early-only contact regardless of attempt budget. Replacing the paired
law by the product adds \((5,\{6\})\), where the only successful slot is
physically unreachable. Adding only \((4,\{7\})\) keeps both marginal
supports but requires two attempts: the two early arrivals cannot tell
which slot works. The late-arrival arm still needs one attempt. The paired
law comes from the fixture, not an empirical statement that cave falls
and upload slot losses correlate. Deleting a slot or retracting the
clearance premise must therefore trigger renewed admission checking,
not validation from the old marginal sets.

\paragraph{Comparison and next cut.}
The \href{https://arxiv.org/html/2103.11470}{NeBula} return-to-network and
data mule give the physical report-delivery actions. The
\href{https://superintelligence.jpl.nasa.gov/public/documents/papers/PlanRob-2020-Planning-and-Optimization-for-Multi-Robot-Planetary-Cave-Exploration.pdf}{JPL
intermittent-connectivity planner} is a route and rendezvous reference.
An \href{https://whoenig.github.io/publications/2019_RA-L_Hoenig.pdf}{action
dependency graph} can express waiting for physical clearance before
departure: the last traversal is a prerequisite for releasing the hold,
not a precedence that can safely be flipped by
\href{https://arxiv.org/abs/2312.04190}{switchable reordering}. This does not
test those implementations or show they cannot be extended. Equipment
authority and an actual local discharge event must be part of any
matched instantiation. Given both, an ordinary timed data-mule controller
with the same history support and a local deadline-aware wait/send
policy game equals the stated budget calculation on 5,616 cases. The 459
positive selected plans also satisfy the route and deadline trace check.
No matched service advantage or independent test clearance is established.
A useful stronger case would let a post-failure diagnostic change both the
lawful hold resource and the reachable return route, comparing a baseline
that may perform the same diagnostic rather than denying it the input.

\end{document}
